% ============================================================
%  中文草稿 —— 请在此版本上直接修改，确认后同步至英文版
%  Chinese Draft: edit here first, then sync to main_en.tex
% ============================================================
\documentclass[12pt]{article}

\usepackage[UTF8]{ctex}           % 中文支持
\usepackage{graphicx}
\usepackage{amsmath, amssymb}
\usepackage{booktabs}
\usepackage{hyperref}
\usepackage{geometry}
\usepackage{cite}
\geometry{a4paper, margin=2.5cm}

% ---------- 论文信息 ----------
\title{\textbf{面向稀疏观测的条件扩散物理一致性\\多视图增强交通速度预测}}
\author{
  严逸锦$^{1}$\quad 邹国健$^{1}$\\[4pt]
  $^{1}$ [机构名称]\\
  \texttt{[邮箱]}
}
\date{2026年3月}

\begin{document}

\maketitle

% ============================================================
\begin{abstract}
% TODO: 摘要 200--250 字
交通速度预测在智能交通系统中扮演着关键角色。然而，现实部署中往往面临
\textbf{训练样本稀缺}与\textbf{节假日分布偏移}两大挑战：
私有路网数据量远少于公开基准（如 PEMS04/08），
且节假日交通模式与工作日分布差异显著。

本文提出 \textbf{DiffAug}——一个基于条件扩散模型的生成式多视图增强框架。
我们训练一个以历���速度–流量–密度 $(v, q, \rho)$ 三元组为条件的 DDPM，
在采样 $K$ 条增强视图后通过物理一致性投影 $\tilde{\rho} = q/(\tilde{v}+\varepsilon)$
严格保证生成速度满足交通流守恒方程。
进一步引入 Beta 插值的 TSMixup 覆盖增强幅度的连续中间态，
并以真实视图为软锚点控制增强权重。

在私有 MSST-CL 数据集（30 个节点，15 分钟粒度）上的实验表明，
DiffAug 相比无增强基线可提升 MAE/RMSE，
同时生成样本的速度分布与真实数据高度一致（KL $< 0.01$，WD $< 0.5$ km/h）。

\textbf{关键词：} 交通速度预测，扩散模型，数据增强，物理约束，稀疏监督
\end{abstract}

% ============================================================
\section{引言}
\label{sec:intro}

% --- 背景与动机 ---
交通速度预测是智能交通系统（ITS）的基础任务之一，
对路径规划、信号控制与拥堵预警均有重要作用。
近年来，以图神经网络（GNN）和 Transformer 为核心的时空预测模型
在公开基准数据集上取得了显著进展~\cite{stgformer,staeformer}。

然而，现实私有路网部署面临两类典型困难：
\begin{enumerate}
  \item \textbf{样本稀缺}：私有路网的历史数据积累有限，
    训练窗口数通常不足公开基准的 1/10，模型容易过拟合。
  \item \textbf{节假日分布偏移}：节假日交通量激增且时段模式与工作日存在系统性偏差，
    模型在节假日上的泛化能力普遍较差。
\end{enumerate}

已有数据增强工作主要分为三类：
（1）\textit{随机结构增强}，如 masking、dropout、时序裁剪，缺乏物理语义约束；
（2）\textit{混合插值}（Mixup~\cite{mixup}），线性插值简单但不保证生成样本的交通物理合法性；
（3）\textit{生成式增强}（STGAN~\cite{stgan}），生成样本与真实分布的一致性难以保证。

本文的核心出发点是：\textit{有效的交通增强视图应当同时满足
（a）分布相似性——与真实交通速度分布高度吻合，
（b）样本多样性——带来工作日训练集中未见的速度场景，
（c）物理合法性——满足 $q = \tilde{\rho} \cdot \tilde{v}$ 的交通流守恒方程。}

基于此，我们提出 \textbf{DiffAug}，核心贡献如下：
\begin{itemize}
  \item 将条件 DDPM 定位为\textbf{数据生成器}而非预测器，
    与预测 backbone 分工解耦，可独立优化；
  \item 设计\textbf{物理一致性投影}模块，在扩散采样后强制满足速度–流量–密度守恒；
  \item 引入\textbf{分布质量筛选}策略，以 KL 散度与 Wasserstein 距离过滤低质量生成视图；
  \item 在私有稀疏路网数据集上验证方案的有效性，
    并通过消融实验分析各模块的贡献。
\end{itemize}

% ============================================================
\section{相关工作}
\label{sec:related}

\subsection{时空交通预测}
% TODO: STGformer, STAEformer, STGCN, DCRNN 等主流方法综述
% 参考 doc/need/相关工作检索与对比模型清单.md

\subsection{交通预测中的数据增强}
% TODO: Mixup, CutMix, STGAN, STAug 等
% 重点区分本文物理约束与已有方法的区别

\subsection{扩散模型用于时序生成}
% TODO: DDPM, DDIM, DiffSTG, TimeDiff 等
% 强调本文"生成器而非预测器"的定位

% ============================================================
\section{问题定义}
\label{sec:problem}

设路网中有 $N$ 个节点，观测周期为 $T$ 步，每步粒度为 15 分钟。
每个时间步的交通状态由速度 $v \in \mathbb{R}^N$、
流量 $q \in \mathbb{R}^N$、密度 $\rho \in \mathbb{R}^N$ 描述，
满足交通流守恒方程：
\begin{equation}
  q = \rho \cdot v.
  \label{eq:traffic_fundamental}
\end{equation}

给定历史观测窗口 $\mathbf{X} = \{(v_t, q_t, \rho_t)\}_{t=1}^{L}$（$L=12$），
目标是预测未来 $H$ 步速度 $\hat{\mathbf{v}} = \{\hat{v}_t\}_{t=L+1}^{L+H}$（$H=12$）。

由于训练样本量受限（约 12,707 个窗口），
我们希望构造 $K$ 个物理合法的增强视图 $\{\tilde{v}^{(k)}\}_{k=1}^K$，
与原始观测共享标签 $\mathbf{Y}$，从而扩充有效训练集规模。

% ============================================================
\section{方法}
\label{sec:method}

\subsection{总体框架}
% TODO: 插入整体架构图

DiffAug 包含三个模块：
\begin{enumerate}
  \item \textbf{条件 DDPM 生成器}（第~\ref{sec:ddpm} 节）：以历史状态为条件采样 $K$ 条速度视图；
  \item \textbf{物理一致性投影}（第~\ref{sec:projection} 节）：将生成速度投影至满足守恒方程的状态空间；
  \item \textbf{TSMixup 插值与加权训练}（第~\ref{sec:mixup} 节）：Beta 插值构造中间增强视图，软锚点控制增强权重。
\end{enumerate}

\subsection{条件扩散生成器}
\label{sec:ddpm}

我们采用 DDPM~\cite{ddpm} 的前向加噪过程：
\begin{equation}
  q(\mathbf{x}_s | \mathbf{x}_0) = \mathcal{N}(\mathbf{x}_s;\, \sqrt{\bar{\alpha}_s}\,\mathbf{x}_0,\, (1-\bar{\alpha}_s)\mathbf{I}),
\end{equation}
以历史速度序列 $\mathbf{x}_0 = v_{1:L}$ 为输入，历史流量、时段编码（ToD）、星期编码（DoW）作为条件 $c$。

去噪网络 $\epsilon_\theta(\mathbf{x}_s, s, c)$ 采用带 ALiBi 位置编码~\cite{alibi} 的 Patch Transformer 结构。
训练目标为标准去噪 MSE：
\begin{equation}
  \mathcal{L}_{\text{denoise}} = \mathbb{E}_{\mathbf{x}_0, \epsilon, s}\left[\|\epsilon - \epsilon_\theta(\mathbf{x}_s, s, c)\|^2\right].
\end{equation}

% TODO: 补充统计损失 (stat_loss) 的公式与动机

\subsection{物理一致性投影}
\label{sec:projection}

DDPM 采样得到增强速度 $\tilde{v}^{(k)}$ 后，密度由以下公式强制满足守恒方程：
\begin{equation}
  \tilde{\rho}^{(k)} = \frac{q}{\tilde{v}^{(k)} + \varepsilon},
  \label{eq:projection}
\end{equation}
其中 $\varepsilon$ 为数值稳定项，$q$ 为原始观测流量（视为已知不变量）。
该投影确保 $(q, \tilde{v}^{(k)}, \tilde{\rho}^{(k)})$ 严格满足式~\eqref{eq:traffic_fundamental}。

\subsection{TSMixup 插值与加权训练}
\label{sec:mixup}

为覆盖原始视图到扩散视图的连续中间态，对每条增强视图做 Beta 插值：
\begin{equation}
  v_{\text{aug}}^{(k)} = \lambda^{(k)}\, v + (1-\lambda^{(k)})\, \tilde{v}^{(k)},
  \quad \lambda^{(k)} \sim \text{Beta}(\alpha, \alpha).
  \label{eq:mixup}
\end{equation}

训练损失由真实视图损失与增强视图损失两部分组成：
\begin{equation}
  \mathcal{L} = \mathcal{L}(v, \mathbf{Y}) + \frac{\alpha_{\text{aug}}}{K}
  \sum_{k=1}^K \mathcal{L}(v_{\text{aug}}^{(k)}, \mathbf{Y}),
  \label{eq:loss}
\end{equation}
其中 $\alpha_{\text{aug}} < 1$ 保证真实观测始终作为监督主轴。

% ============================================================
\section{实验}
\label{sec:exp}

\subsection{数据集与实验设置}

\paragraph{MSST-CL 数据集}
% TODO: 数据集描述
私有路网，$N=30$ 个节点，时间粒度 15 分钟，
训练/验证/测试划分为 60\%/20\%/20\%，
训练集约 12,707 个历史窗口（$L=12$，$H=12$）。
包含节假日段，存在显著分布偏移。

\paragraph{基线方法}
% TODO: 列出对比方法
STGformer~\cite{stgformer}、
STGCN~\cite{stgcn}、
DCRNN~\cite{dcrnn}、
STAEformer~\cite{staeformer}、
DiffSTG~\cite{diffstg}。

\paragraph{评价指标}
MAE、RMSE、MAPE（$H=12$ 步平均）。

\subsection{主实验结果}
% TODO: 插入主实验表格

\subsection{消融实验}

表~\ref{tab:ablation} 展示了各模块的贡献：
\begin{table}[ht]
  \centering
  \caption{消融实验结果（MSST-CL 数据集，MAE/RMSE）}
  \label{tab:ablation}
  \begin{tabular}{lcccc}
    \toprule
    \textbf{Variant} & \textbf{策略} & \textbf{KL}$\downarrow$ & \textbf{MAE}$\downarrow$ & \textbf{RMSE}$\downarrow$ \\
    \midrule
    no\_aug       & —                & —      & 4.4045 & 6.6519 \\
    e4            & stat\_loss       & 0.0022 & 4.4225 & 6.6494 \\
    e1            & adaLN            & —      & 4.4473 & 6.6907 \\
    baseline\_aug & —                & —      & 4.4413 & 6.6859 \\
    \midrule
    \textbf{DiffAug (ours)} & e3e4 + 筛选 & \textbf{TODO} & \textbf{TODO} & \textbf{TODO} \\
    \bottomrule
  \end{tabular}
\end{table}

\subsection{生成质量分析}
% TODO: 插入 KL、WD、intra_std 分析图/表

% ============================================================
\section{结论}
\label{sec:conclusion}

本文提出 DiffAug，一个面向稀疏观测交通路网的条件扩散生成式多视图增强框架。
通过物理一致性投影与 TSMixup 插值，
DiffAug 在保证生成速度分布高度吻合真实交通分布的同时，
引入了工作日训练集中未见的速度多样性。
实验结果表明，DiffAug 在私有稀疏路网数据集上能够有效缓解过拟合，
在节假日场景下泛化能力显著优于无增强基线。

\paragraph{局限性与未来工作}
当前筛选策略依赖离线预计算，未来可探索在线自适应筛选；
此外，将扩散生成与预测 backbone 端到端联合训练也值得进一步研究。

% ============================================================
\bibliographystyle{plain}
\bibliography{refs}

% ============================================================
\appendix
\section{附录：扩散模型 Variant 详细配置}
% TODO: 各 variant 的超参数表

\end{document}
